% The forecast timeline, from figures/forecast-timeline.mmd.
%
% The diagram itself is Mermaid; this file only places it and carries the
% caption and label, so \ref{fig:forecast-timeline} resolves the same way it
% did when the figure was drawn in TikZ.
%
% To regenerate after editing the .mmd:
%   reporting/regenerate.sh timeline
% or by hand, from this folder:
%   npx -y @mermaid-js/mermaid-cli -i forecast-timeline.mmd \
%       -o ../build/figures/forecast-timeline.pdf -b white --pdfFit
% --pdfFit trims the page to the diagram. The output is vector, so it stays
% sharp at any size; a PNG of this needs -s 4 (a scale factor -- mermaid sizes
% the canvas to the diagram and ignores -w) just to print acceptably.
%
% The TikZ original is still here as forecast-timeline.tex. To go back, point
% content.tex at it instead of at this file -- nothing else changes.
%
% The .pdf is found through \graphicspath in main.tex, which looks in
% ../build/figures. Naming it without a path is what lets this wrapper be
% \input from anywhere.
\begin{figure}[htbp]
\centering
% Width knob: the .mmd sets the diagram's own proportions, this sets how
% much of the text block it takes. Raise or lower the 0.85 to taste.
\includegraphics[width=0.85\textwidth]{forecast-timeline.pdf}
\caption[How each model reports a forecast day]{%
  \textbf{What each model means by a forecast day.}
  One column per forecast day, one row per model.
  A \textbf{solid} block names the days it covers, so it earns those days a
  place on the timeline. A \textbf{dashed} block covers days without naming
  any: it is drawn on each of them, but cannot by itself make a day worth
  stepping to. The Flood Hub block sits on days 5--6 purely as an
  illustration: a published window is one to two days wide and can fall
  anywhere in the first week, and which days it names depends on the gauge.
  Colour here separates the models in the diagram only --- on the map they
  share a palette and are told apart by the panel.}
\label{fig:forecast-timeline}
\end{figure}
